\subsection{Адресация: где какой светодиод?}

Мы выяснили, что в буфере 25 ячеек (\texttt{0..24}). Но на плате квадрокоптера вы, скорее всего, видите только 4 больших светодиода (или 29, если установлен модуль расширения).

Как понять, какой индекс в коде соответствует какой лампочке в реальности?

\begin{taskbox}[Исследовательская задача]
Напишите программу-исследователь, которая поможет составить «карту» светодиодов.
\begin{enumerate}
    \item Включите светодиод с индексом \texttt{0} красным цветом.
    \item Посмотрите на дрон. Какой светодиод загорелся? (Например, «Передний левый»).
    \item Повторите для индексов \texttt{1}, \texttt{2}, \texttt{3}.
    \item Запишите результаты в тетрадь.
\end{enumerate}
\end{taskbox}

\begin{codebox}[Скрипт-исследователь]
\begin{minted}{lua}
local leds = Ledbar.new(25)

-- Функция для очистки (выключения всех)
local function clear()
    for i=0, 24 do leds:set(i, 0,0,0) end
end

clear()
-- Включаем индекс 0
leds:set(0, 1, 0, 0) 
-- Если загорелся не тот, который вы ожидали — поменяйте 0 на 1, 2 или 3
\end{minted}
\end{codebox}

Это упражнение учит главному принципу инженера: \textbf{если документации недостаточно — проведи эксперимент!}
